\documentclass{article}


% if you need to pass options to natbib, use, e.g.:
%     \PassOptionsToPackage{numbers, compress}{natbib}
% before loading neurips_2023


% ready for submission
\usepackage{neurips_2023}
\input{math_commands}

% to compile a preprint version, e.g., for submission to arXiv, add add the
% [preprint] option:
%     \usepackage[preprint]{neurips_2023}


% to compile a camera-ready version, add the [final] option, e.g.:
%     \usepackage[final]{neurips_2023}


% to avoid loading the natbib package, add option nonatbib:
%    \usepackage[nonatbib]{neurips_2023}


\usepackage[utf8]{inputenc} % allow utf-8 input
\usepackage[T1]{fontenc}    % use 8-bit T1 fonts
\usepackage{hyperref}       % hyperlinks
\usepackage{url}            % simple URL typesetting
\usepackage{booktabs}       % professional-quality tables
\usepackage{amsfonts}       % blackboard math symbols
\usepackage{nicefrac}       % compact symbols for 1/2, etc.
\usepackage{microtype}      % microtypographyhttps://www.overleaf.com/project/64c80f60d7cc8f8216a032d6
\usepackage{xcolor}         % colors
\usepackage{graphicx}
\usepackage{subcaption}



\title{Taking Features out of Superposition with Sparse Dictionary Learning}


% The \author macro works with any number of authors. There are two commands
% used to separate the names and addresses of multiple authors: \And and \AND.
%
% Using \And between authors leaves it to LaTeX to determine where to break the
% lines. Using \AND forces a line break at that point. So, if LaTeX puts 3 of 4
% authors names on the first line, and the last on the second line, try using
% \AND instead of \And before the third author name.


\author{Hoagy Cunningham\thanks{Equal contribution}\hspace{4pt}$^{1 2}$,\  Aidan Ewart$^{* 1 3}$, \ Logan Riggs$^{* 1}$, \ Robert Huben, \ Lee Sharkey$^{4}$ \\
$^1$EleutherAI, $^2$MATS, $^3$Bristol AI Safety Centre, $^4$Apollo Research \\
\texttt{\{hoagycunningham, aidanprattewart, logansmith5\}@gmail.com}}


\begin{document}


\maketitle


\begin{abstract}
One of the roadblocks to a better understanding of neural networks' internals is \textit{polysemanticity}, where neurons appear to activate in multiple, semantically distinct contexts. Polysemanticity prevents us from identifying concise, human-understandable explanations for what neural networks are doing internally. One hypothesised cause of polysemanticity is \textit{superposition}, where neural networks represent more features than they have neurons by assigning features to an overcomplete set of directions in activation space, rather than to individual neurons. Here, we attempt to identify those directions, using sparse autoencoders to reconstruct the internal activations of a language model. These autoencoders learn sets of sparsely activating features that are more interpretable and monosemantic than directions identified by alternative approaches, where interpretability is measured by automated methods. This work indicates that it is possible to resolve superposition in language models using a scalable, unsupervised method. Our method may serve as a foundation for future mechanistic interpretability work, which we hope will enable greater model transparency and steerability.

\end{abstract}


\section{Introduction}
\input{neurips_workshop/introduction_neurips}

\section{Taking Features out of Superposition with Sparse Dictionary Learning}
\label{section:dictionary_learning}
\input{neurips_workshop/dictionary_learning_neurips}

\section{Interpreting Dictionary Features}
\label{section:autointerp}
\subsection{Interpretability at Scale}
\label{section:autointerp_background}
\input{sections/autointerpretability_background}

\subsection{Sparse dictionary features are more interpretable than baselines}
\label{subsection:interpretability_baselines}
\input{neurips_workshop/autointerp_neurips}

\iffalse
\section{Dictionary Features Enable Surgical Concept Erasure}

\label{section:concept_erasure}

\input{sections/aidan_concept_editing_results}
\fi

\section{Case Studies}

\label{section:case_studies}

\input{sections/logan_case_studies}


\section{Discussion}
\input{neurips_workshop/neurips_conclusion}


% \subsubsection*{Acknowledgments}
% \input{sections/acknowledgements}


\bibliography{iclr2024_conference_google_scholar_version}
\bibliographystyle{iclr2024_conference}

\appendix

\section{Autointerpretation Protocol}
\label{appendix:autointerp_protocol}
\input{sections/appendices/autointerp_method_details}

\section{Sparse Autoencoder Training and Hyperparameter Selection}
\label{appendix:hparam_selection}
\input{sections/appendices/hyperparameter_selection}

\section{Further Autointerpretation Results}
\label{appendix:full_autointerp}
\input{neurips_workshop/further_autointerpretability_results_neurips}

\iffalse
\section{Measuring the Disruptiveness of our Concept Erasure}
\label{appendix:causal_interventions}
\input{sections/appendices/aidan_erased_feature_interp}
\fi

\section{Identifying Causally-Important Dictionary Features for Indirect Object Identification}
\label{section:identifying_causal_features}
\input{sections/aidan_identifying_causal_features}

\section{Qualitative Feature Analysis}
\label{appendix:feature_interp_examples}
\input{sections/appendices/feature_analysis}

\section{Number of Active Features}
\label{appendix:n_features_active}
\input{sections/appendices/prop_feats_active}

% \section{Outlier Dimensions}
% \label{appendix:outlier_dimensions}
% \input{sections/appendices/outlier_dimensions}

% \section{Sparse Autoencoding non-residual-stream stuff}



\section{Editing IOI Behaviour on other Layers}
\label{appendix:ioi_other_layers}
\input{sections/appendices/aidan_ioi_other_layers}

\section{Top K Comparisons}
\label{appendix:topk}
\input{sections/appendices/full_baselines}


% \subsubsection*{Author Contributions}
% If you'd like to, you may include  a section for author contributions as is done
% in many journals. This is optional and at the discretion of the authors.




\end{document}
